\section{PHÂN TÍCH VÀ KẾT LUẬN}

\subsection{Phân tích độ trễ và khả năng thông lượng (Throughput)}
Thông qua công cụ \texttt{performance\_test.py}, lưu lượng trao đổi trên hệ thống được đo đạc lặp lại ở hai phân vùng chính: Mạng cục bộ không qua Tường lửa (Inside-to-Inside) và Mạng kết nối biên vượt qua Tường lửa (Inside-to-DMZ, External-to-DMZ). Mọi kết quả sau khi thu thập được tự động kết xuất ra một tệp báo cáo Excel, có đính kèm biểu đồ Matplotlib thông qua thư viện \texttt{openpyxl}.

% CHỤP ẢNH MINH CHỨNG:
% Yêu cầu: Mở file báo cáo Excel (ví dụ: results/mmtnc_report_20260423_214502.xlsx), mở Sheet "Throughput", chụp màn hình nội dung Sheet này (cả bảng số liệu và biểu đồ).
% Lưu ảnh vào thư mục image/ với tên report_excel.png và bỏ comment đoạn code bên dưới để hiển thị ảnh.
% \begin{figure}[H]
%     \centering
%     \includegraphics[width=0.9\textwidth]{image/report_excel.png}
%     \caption{Kết xuất báo cáo đo lường thông lượng qua tệp Excel tự động}
%     \label{fig:report_excel}
% \end{figure}

Dữ liệu cho thấy:
\begin{itemize}
    \item \textbf{Độ trễ trung bình (RTT Latency):} Luồng dữ liệu đi qua biên mạng (phải xử lý chuyển đổi địa chỉ NAT, và lọc qua ACL) có độ trễ lớn hơn khi so sánh với định tuyến tĩnh trực tiếp trong lớp nội bộ. Nguyên nhân xuất phát từ chi phí trích xuất gói tin của Tường lửa.
    \item \textbf{Thông lượng truyền tải TCP:} Lượng băng thông cao nhất đo được qua iPerf cho thấy lớp Distribution và Core hoạt động hiệu quả đạt được tốc độ cận tối đa do cấu hình đường truyền gigabit ảo. 
\end{itemize}

\subsection{Đánh giá tác động của tiến trình xử lý NAT (NAT Overhead)}
Hệ thống tiến hành thực hiện một bài toán so sánh độ trễ truyền gói tin giữa kịch bản máy tính có khả năng ping trực tiếp mạng ngoài (Route thuần túy) với kịch bản máy tính bị thay đổi cấu trúc gói tin (MASQUERADE).
Qua tập dữ liệu \texttt{NAT\_Overhead} trên báo cáo tự động, ta nhận thấy việc bổ sung NAT làm tăng thời gian RTT (Round Trip Time). Do chi phí phải tra cứu bảng trạng thái (Connection Tracking) của iptables, hiệu năng định tuyến giảm tương đối. Tuy nhiên, mức độ trễ (delay) rơi vào vài phần ngàn giây (milliseconds), hoàn toàn ở mức chấp nhận được để đánh đổi lại lớp rào cản phòng vệ kiên cố cũng như tiết kiệm IP Public.

\subsection{Kết luận}
Thông qua việc vận dụng các cấu trúc định tuyến tĩnh và động trong môi trường Mininet kết hợp FRRouting, đồ án đã triển khai thành công một mô hình \textbf{Campus 3-tier tích hợp DMZ}. Việc quy hoạch địa chỉ rõ ràng kết hợp cùng các giao thức tự động hóa đã đáp ứng xuất sắc các mục tiêu sau:
\begin{enumerate}
    \item Tách biệt vùng mạng nhạy cảm (Inside) với vùng dịch vụ kết nối ngoài (DMZ).
    \item Sử dụng iptables (Tường lửa) để bảo vệ toàn diện ứng dụng dịch vụ trên máy chủ (qua cơ chế ACL và NAT) song song với quản trị nguồn IP (PAT).
    \item Áp dụng thành công tiến trình OSPF giúp hệ thống dễ dàng mở rộng và định tuyến chính xác.
    \item Tính năng định tuyến tự động khi xảy ra tắc nghẽn (Cân bằng tải theo ngưỡng) góp phần hiện thực hoá nguyên lý sẵn sàng và giảm thiểu thời gian chết của dịch vụ (Downtime).
\end{enumerate}

Những đánh giá về mặt hiệu suất là minh chứng rõ rệt cho chi phí về thời gian xử lý khi thực hiện bộ đệm tường lửa. Mô hình hoàn toàn có khả năng mở rộng để ứng dụng các kịch bản mạng máy tính phức tạp hơn trong tương lai, cung cấp một nền tảng thực tiễn về cấu trúc an ninh hệ thống mạng hiện đại.
